
\documentclass{article}
\usepackage{amsfonts}

%%%%%%%%%%%%%%%%%%%%%%%%%%%%%%%%%%%%%%%%%%%%%%%%%%%%%%%%%%%%%%%%%%%%%%%%%%%%%%%%%%%%%%%%%%%%%%%%%%%%%%%%%%%%%%%%%%%%%%%%%%%%%%%%%%%%%%%%%%%%%%%%%%%%%%%%%%%%%%%%%%%%%%%%%%%%%%%%%%%%%%%%%%%%%%%%%%%%%%%%%%%%%%%%%%%%%%%%%%%%%%%%%%%
\usepackage{amsmath, amssymb}
\usepackage{geometry}

\setcounter{MaxMatrixCols}{10}
%TCIDATA{OutputFilter=LATEX.DLL}
%TCIDATA{Version=5.50.0.2960}
%TCIDATA{<META NAME="SaveForMode" CONTENT="1">}
%TCIDATA{BibliographyScheme=Manual}
%TCIDATA{LastRevised=Monday, March 10, 2025 07:26:38}
%TCIDATA{<META NAME="GraphicsSave" CONTENT="32">}

\geometry{a4paper, margin=1in}

\input{tcilatex}

\begin{document}

\title{DeepSeek Answers on Clifford Algebras and Dirac Matrices}
\author{}
\date{}
\maketitle

\section*{Clifford Algebra $\text{Cl}(4)$ and Tensor Products}

The Clifford algebra $\text{Cl}(2)$ is a 4-dimensional algebra generated by
two elements $e_1$ and $e_2$ that satisfy the relations: 
\begin{equation*}
e_1^2 = e_2^2 = 1, \quad e_1 e_2 = -e_2 e_1. 
\end{equation*}
The basis of $\text{Cl}(2)$ is $\{1, e_1, e_2, e_1 e_2\}$, and it is
isomorphic to the algebra of $2 \times 2$ matrices over the real numbers, $%
\text{Mat}(2, \mathbb{R})$.

The tensor product $\text{Cl}(2) \otimes \text{Cl}(2)$ is isomorphic to $%
\text{Cl}(4)$, the Clifford algebra generated by four elements $f_1, f_2,
f_3, f_4$ satisfying: 
\begin{equation*}
f_i^2 = 1 \quad \text{for all } i, \quad f_i f_j = -f_j f_i \quad \text{for }
i \neq j. 
\end{equation*}
The basis of $\text{Cl}(4)$ consists of $2^4 = 16$ elements, which are all
possible products of the generators $f_1, f_2, f_3, f_4$.

\section*{Map of $\text{Cl}(2)$ into $\text{Cl}(4)$}

The map $\phi: \text{Cl}(2) \to \text{Cl}(4)$ can be defined as: 
\begin{equation*}
\phi(1) = 1, \quad \phi(e_1) = f_1, \quad \phi(e_2) = f_2, \quad \phi(e_1
e_2) = f_1 f_2. 
\end{equation*}
This map preserves the algebraic structure of $\text{Cl}(2)$ because the
relations in $\text{Cl}(2)$ are satisfied in $\text{Cl}(4)$: 
\begin{equation*}
f_1^2 = 1, \quad f_2^2 = 1, \quad f_1 f_2 = -f_2 f_1. 
\end{equation*}

\section*{Skew Tensor Product}

The skew tensor product (graded tensor product) of two graded algebras $A$
and $B$ is defined as: 
\begin{equation*}
(a \otimes b) \cdot (a^{\prime }\otimes b^{\prime |b| \cdot |a^{\prime }|}
(a a^{\prime }) \otimes (b b^{\prime }), 
\end{equation*}
where $|a|$ and $|b|$ are the degrees of $a$ and $b$, respectively. This
ensures that the multiplication respects the grading and anticommutation
relations.

\section*{Dirac Matrices and $\text{Cl}(1, 3)$}

The Clifford algebra $\text{Cl}(1, 3)$ is generated by the Dirac matrices $%
\gamma_\mu$ ($\mu = 0, 1, 2, 3$) satisfying: 
\begin{equation*}
\{ \gamma_\mu, \gamma_\nu \} = 2 \eta_{\mu\nu} I, 
\end{equation*}
where $\eta_{\mu\nu}$ is the Minkowski metric with signature $(+, -, -, -)$.
The Dirac matrices act on 4-component Dirac spinors, which form a
4-dimensional representation of the Lorentz group.

\section*{Complex Clifford Algebras}

In the complex case, the Clifford algebra $\text{Cl}(p, q)$ depends only on
the total number of generators $n = p + q$, not on the signature $(p, q)$.
For example: 
\begin{equation*}
\text{Cl}(1, 3) \otimes \mathbb{C} \cong \text{Cl}(4, 0) \otimes \mathbb{C}
\cong \text{Cl}(0, 4) \otimes \mathbb{C} \cong \text{Mat}(4, \mathbb{C}). 
\end{equation*}
This is because the complex numbers allow us to rescale generators to
eliminate the distinction between timelike and spacelike directions.

\end{document}
